% answers_short.tex —— 1985 年·上海卷·理工农医类 · 简短答案速查
% 由 tools/make_short_answers.py 自 answers.tex 的 \gksolhead/\gkanswer 提取，
% 并与 1985/1985_数学真题_上海理_答案要点.md 逐题比对后生成，勿手改（改源 MD 后重跑）。

\begin{gknotice}
\textbf{使用说明}：本节只给最终答案便于快速核对，完整推导见解析版。答案为编写者独立推导，\textbf{非官方标准答案}。
\end{gknotice}

\gkpart{大题}
\begin{itemize}[leftmargin=2.6em]
  \item[\textbf{1.}] (1) $(-3,-1)$；(2) $y=(x+1)^8\ (x\ge-1)$；(3) $\dfrac{\sqrt2}{2}$；(4) $4$；(5) $(x-2)^2+(y-4)^2=5$；(6) $l\parallel\alpha$；(7) $-1$；(8) $x=2$；(9) $\dfrac{\sqrt5}{5}$；(10) $\pi$
  \item[\textbf{2.}] (1) （C）；(2) （C）；(3) （D）；(4) （B）；(5) （D）
  \item[\textbf{3.}] (1) 双曲线 $\dfrac{y^2}{4}-\dfrac{x^2}{9}=1$；(2) 圆 $x^2+(y-1)^2=1$
  \item[\textbf{4.}] (1) $-5005$；(2) 15 种取法，180 个四位偶数；(3) $\dfrac{16\pi}{3}$
  \item[\textbf{5.}] $z_1+z_2$ 为实数。
  \item[\textbf{6.}] $x=\dfrac\pi4+k\dfrac\pi2$ 或 $x=\dfrac\pi8+k\dfrac\pi4$，$k\in\mathbb Z$。
  \item[\textbf{7.}] 方程 $f(x)=0$ 至多有一个实数根。
  \item[\textbf{8.}] 按数学归纳法证明题面不等式成立。
  \item[\textbf{9.}] $\alpha_1=\alpha_2=\alpha_3$。
  \item[\textbf{10.}] 见解析正文。
\end{itemize}
